\chapter{Differentiable Manifolds}
\label{chap:dc-differentiable-manifolds}

The model example in \cite{doCarmo} is a regular surface $S\subset\mathbb{R}^3$:
for every $p\in S$, some open $U\subset\mathbb{R}^2$ and neighborhood
$V\subset\mathbb{R}^3$ admit a map
$\mathbf{x}:U\to V\cap S$ that is a differentiable homeomorphism onto $V\cap S$ and
has injective differential at every point. For overlapping parametrizations, the
transition maps are diffeomorphisms (\cref{ex:dc-ch0-4-2}). The abstract definition
below replaces the ambient $\mathbb{R}^3$ by $\mathbb{R}^n$.

Throughout, ``differentiable'' means $C^\infty$.

\section{Differentiable manifolds; tangent space}

\begin{definition}[differentiable manifold]
  \label{def:dc-ch0-2-1}
  \lean{Riemannian.DCDifferentiableManifold}
  \leanok
  \dcref{ch0:2.1}
  \source{docarmo:page-0017}
  A \emph{differentiable manifold} of dimension $n$ is a set $M$ together with a family
  of injective mappings $\mathbf{x}_\alpha:U_\alpha\subset\mathbb{R}^n\to M$ of open
  sets $U_\alpha$ of $\mathbb{R}^n$ into $M$ such that:
  1. $\bigcup_\alpha\mathbf{x}_\alpha(U_\alpha)=M$.
  2. for any pair $\alpha,\beta$ with
     $\mathbf{x}_\alpha(U_\alpha)\cap\mathbf{x}_\beta(U_\beta)=W\neq\phi$, the sets
     $\mathbf{x}_\alpha^{-1}(W)$ and $\mathbf{x}_\beta^{-1}(W)$ are open in
     $\mathbb{R}^n$ and the mappings $\mathbf{x}_\beta^{-1}\circ\mathbf{x}_\alpha$
     are differentiable.
  3. the family $\{(U_\alpha,\mathbf{x}_\alpha)\}$ is maximal relative to (1) and (2).
  The pair $(U_\alpha,\mathbf{x}_\alpha)$ (or $\mathbf{x}_\alpha$) with
  $p\in\mathbf{x}_\alpha(U_\alpha)$ is a \emph{parametrization} (or \emph{system of
  coordinates}) of $M$ at $p$; $\mathbf{x}_\alpha(U_\alpha)$ is then a \emph{coordinate
  neighborhood} at $p$. A family satisfying (1) and (2) is a \emph{differentiable
  structure} on $M$. Any structure satisfying (1)--(2) has a unique maximal extension.
\end{definition}

\begin{remark}[change of parameters as an axiom]
  \label{rem:dc-ch0-2-2}
  \uses{def:dc-ch0-2-1}
  \notready
  \dcref{ch0:2.2}
  \notready
  Condition (2) of \cref{def:dc-ch0-2-1} declares every coordinate change
  differentiable; it is the
  compatibility condition for transferring calculus from $\mathbb{R}^n$ to $M$.
\end{remark}

\begin{remark}[natural topology]
  \label{rem:dc-ch0-2-3}
  \notready
  \dcref{ch0:2.3}
  \notready
  A differentiable structure on $M$ induces a natural topology: define $A\subset M$
  to be \emph{open} if and only if $\mathbf{x}_\alpha^{-1}(A\cap\mathbf{x}_\alpha(U_\alpha))$
  is open in $\mathbb{R}^n$ for all $\alpha$. This defines a topology: $M$ and $\phi$
  are open, arbitrary unions and finite intersections are open, each coordinate
  neighborhood is open, and every $\mathbf{x}_\alpha$ is continuous. The identity
  atlas gives the usual differentiable structure on $\mathbb{R}^n$.
\end{remark}

\begin{example}[the real projective space $P^n(\mathbb{R})$]
  \label{ex:dc-ch0-2-4}
  \notready
  \dcref{ch0:2.4}
  \notready
  $P^n(\mathbb{R})$ is the set of straight lines of $\mathbb{R}^{n+1}$ through the
  origin, i.e. the quotient of $\mathbb{R}^{n+1}-\{0\}$ by
  $(x_1,\dots,x_{n+1})\sim(\lambda x_1,\dots,\lambda x_{n+1})$, $\lambda\neq0$,
  with points denoted $[x_1,\dots,x_{n+1}]$. Let
  $V_i=\{[x_1,\dots,x_{n+1}]; x_i\neq0\}$, $i=1,\dots,n+1$, and define
  $\mathbf{x}_i:\mathbb{R}^n\to V_i$ by
  $$
  \mathbf{x}_i(y_1,\dots,y_n)=[y_1,\dots,y_{i-1},1,y_i,\dots,y_n].
  $$
  Each $\mathbf{x}_i$ is bijective with $\bigcup\mathbf{x}_i(\mathbb{R}^n)=P^n(\mathbb{R})$,
  $\mathbf{x}_i^{-1}(V_i\cap V_j)$ is open, and (for $i>j$, the case $i<j$ similar)
  $$
  \mathbf{x}_j^{-1}\circ\mathbf{x}_i(y_1,\dots,y_n)=\Big(\frac{y_1}{y_j},\dots,\frac{y_{j-1}}{y_j},\frac{y_{j+1}}{y_j},\dots,\frac{y_{i-1}}{y_j},\frac{1}{y_j},\frac{y_i}{y_j},\dots,\frac{y_n}{y_j}\Big)
  $$
  is differentiable. Thus $\{(\mathbb{R}^n,\mathbf{x}_i)\}$ is a differentiable
  structure: $P^n(\mathbb{R})$ is covered by $n+1$ coordinate neighborhoods $V_i$
  (directions not in the hyperplane $x_i=0$), with \emph{inhomogeneous coordinates}
  $\big(\frac{x_1}{x_i},\dots,\frac{x_{i-1}}{x_i},\frac{x_{i+1}}{x_i},\dots,\frac{x_{n+1}}{x_i}\big)$
  corresponding to the \emph{homogeneous coordinates} $(x_1,\dots,x_{n+1})$.
\end{example}

\begin{definition}[differentiable mapping]
  \label{def:dc-ch0-2-5}
  \lean{Riemannian.DCDifferentiableAt}
  \uses{def:dc-ch0-2-1}
  \leanok
  \dcref{ch0:2.5}
  Let $M_1^n$ and $M_2^m$ be differentiable manifolds. A mapping
  $\varphi:M_1\to M_2$ is \emph{differentiable at} $p\in M_1$ if, given a parametrization
  $\mathbf{y}:V\subset\mathbb{R}^m\to M_2$ at $\varphi(p)$, there is a
  parametrization $\mathbf{x}:U\subset\mathbb{R}^n\to M_1$ at $p$ with
  $\varphi(\mathbf{x}(U))\subset\mathbf{y}(V)$ such that
  $$
  \mathbf{y}^{-1}\circ\varphi\circ\mathbf{x}:U\subset\mathbb{R}^n\to\mathbb{R}^m\qquad\text{(1)}
  $$
  is differentiable at $\mathbf{x}^{-1}(p)$. By condition (2) of \cref{def:dc-ch0-2-1}
  this is independent of the choice of parametrizations. The mapping (1) is the
  \emph{expression} of $\varphi$ in the parametrizations $\mathbf{x}$ and $\mathbf{y}$.
\end{definition}

\begin{definition}[tangent vector; tangent space]
  \label{def:dc-ch0-2-6}
  \lean{Riemannian.DCTangentSpaceAt}
  \leanok
  \dcref{ch0:2.6}
  Let $M$ be a differentiable manifold. A differentiable
  $\alpha:(-\varepsilon,\varepsilon)\to M$ is a \emph{curve} in $M$. Suppose
  $\alpha(0)=p$, and let $\mathcal{D}$ be the functions on $M$ differentiable at
  $p$. The \emph{tangent vector to the curve $\alpha$ at $t=0$} is the function
  $\alpha'(0):\mathcal{D}\to\mathbb{R}$ given by
  $$
  \alpha'(0)f=\frac{d(f\circ\alpha)}{dt}\Big|_{t=0},\qquad f\in\mathcal{D}.
  $$
  A \emph{tangent vector at $p$} is the tangent vector at $t=0$ of some curve
  $\alpha:(-\varepsilon,\varepsilon)\to M$ with $\alpha(0)=p$. The set of all such
  is denoted $T_pM$. Choosing a parametrization $\mathbf{x}:U\to M^n$ at
  $p=\mathbf{x}(0)$, with $\mathbf{x}^{-1}\circ\alpha(t)=(x_1(t),\dots,x_n(t))$,
  one obtains
  $$
  \alpha'(0)=\sum_i x_i'(0)\Big(\frac{\partial}{\partial x_i}\Big)_0.\qquad\text{(2)}
  $$
  Here $\big(\frac{\partial}{\partial x_i}\big)_0$ is the tangent vector at $p$ of
  the coordinate curve $x_i\mapsto\mathbf{x}(0,\dots,x_i,\dots,0)$. Expression (2)
  shows $T_pM$ is a vector space of dimension $n$ with \emph{associated basis}
  $\big\{\big(\frac{\partial}{\partial x_1}\big)_0,\dots,\big(\frac{\partial}{\partial x_n}\big)_0\big\}$,
  independent of the parametrization; $T_pM$ is the \emph{tangent space} of $M$ at $p$.
\end{definition}

\begin{proposition}[the differential]
  \label{prop:dc-ch0-2-7}
  \lean{Riemannian.DCDifferentialAt_curve}
  \uses{def:dc-ch0-2-6}
  \leanok
  \dcref{ch0:2.7}
  Let $M_1^n,M_2^m$ be differentiable manifolds and $\varphi:M_1\to M_2$
  differentiable. For $p\in M_1$ and $v\in T_pM_1$, choose a differentiable curve
  $\alpha:(-\varepsilon,\varepsilon)\to M_1$ with $\alpha(0)=p$, $\alpha'(0)=v$, and
  put $\beta=\varphi\circ\alpha$. Then
  $d\varphi_p:T_pM_1\to T_{\varphi(p)}M_2$ given by $d\varphi_p(v)=\beta'(0)$ is a
  linear mapping independent of the chosen curve $\alpha$.
\end{proposition}
\begin{proof}
  \sketch
  \leanok
  In parametrizations $\mathbf{x}$ at $p$ and $\mathbf{y}$ at $\varphi(p)$,
  write $\mathbf{y}^{-1}\circ\varphi\circ\mathbf{x}(q)=(y_1(x_1,\dots,x_n),\dots,y_m(x_1,\dots,x_n))$
  and $\mathbf{x}^{-1}\circ\alpha(t)=(x_1(t),\dots,x_n(t))$. Then with respect to
  the basis $\{(\frac{\partial}{\partial y_i})_0\}$,

  $$
  \beta'(0)=\Big(\sum_{i=1}^n\frac{\partial y_1}{\partial x_i}x_i'(0),\dots,\sum_{i=1}^n\frac{\partial y_m}{\partial x_i}x_i'(0)\Big),\qquad\text{(3)}
  $$
  
  which is independent of $\alpha$ and equals
  $d\varphi_p(v)=\big(\frac{\partial y_i}{\partial x_j}\big)(x_j'(0))$. Thus
  $d\varphi_p$ is linear with matrix $\big(\frac{\partial y_i}{\partial x_j}\big)$
  in the associated bases.
\end{proof}

\begin{definition}[differential]
  \label{def:dc-ch0-2-8}
  \lean{Riemannian.DCDifferentialAt}
  \uses{prop:dc-ch0-2-7}
  \leanok
  \dcref{ch0:2.8}
  The linear mapping $d\varphi_p$ defined by \cref{prop:dc-ch0-2-7} is the \emph{differential}
  of $\varphi$ at $p$.
\end{definition}

\begin{definition}[diffeomorphism; local diffeomorphism]
  \label{def:dc-ch0-2-9}
  \lean{Riemannian.DCDiffeomorph}
  \uses{def:dc-ch0-2-5}
  \leanok
  \dcref{ch0:2.9}
  $\varphi:M_1\to M_2$ is a \emph{diffeomorphism} if it is differentiable, bijective, and
  $\varphi^{-1}$ is differentiable. It is a \emph{local diffeomorphism at} $p$ if there
  exist neighborhoods $U$ of $p$ and $V$ of $\varphi(p)$ with $\varphi:U\to V$ a
  diffeomorphism. By the chain rule, if $\varphi$ is a diffeomorphism then
  $d\varphi_p:T_pM_1\to T_{\varphi(p)}M_2$ is an isomorphism for all $p$; in
  particular $\dim M_1=\dim M_2$.
\end{definition}

\begin{lemma}[differential of a diffeomorphism is an isomorphism]
  \label{lem:dc-ch0-2-9-iso}
  \lean{Riemannian.DCDiffeomorph.mfderiv_bijective}
  \uses{def:dc-ch0-2-9, def:dc-ch0-2-8}
  \leanok
  \dcref{ch0:2.9}
  If $\varphi:M_1\to M_2$ is a diffeomorphism then for every $p\in M_1$ the
  differential $d\varphi_p:T_pM_1\to T_{\varphi(p)}M_2$ is a linear isomorphism.
\end{lemma}
\begin{proof}
  \sketch
  \leanok
  Let $\psi=\varphi^{-1}$. By the chain rule (\cref{prop:dc-ch0-2-7}) applied to
  $\psi\circ\varphi=\mathrm{id}_{M_1}$ and $\varphi\circ\psi=\mathrm{id}_{M_2}$,
  $d\psi_{\varphi(p)}\circ d\varphi_p=\mathrm{id}_{T_pM_1}$ and
  $d\varphi_p\circ d\psi_{\varphi(p)}=\mathrm{id}_{T_{\varphi(p)}M_2}$. Hence
  $d\varphi_p$ is a two-sided invertible linear map, i.e. a linear isomorphism, with
  inverse $d\psi_{\varphi(p)}$.
\end{proof}

\begin{theorem}[local converse, via inverse function theorem]
  \label{thm:dc-ch0-2-10}
  \lean{Riemannian.isLocalDiffeomorphAt_of_mfderiv_equiv}
  \uses{def:dc-ch0-2-5, def:dc-ch0-2-8, def:dc-ch0-2-9}
  \leanok
  \dcref{ch0:2.10}
  Let $\varphi:M_1^n\to M_2^n$ be a differentiable mapping and $p\in M_1$ such that
  $d\varphi_p:T_pM_1\to T_{\varphi(p)}M_2$ is an isomorphism. Then $\varphi$ is a
  local diffeomorphism at $p$.
\end{theorem}
\begin{proof}
  \sketch
  \leanok
  Choose coordinate charts $\mathbf{x}_1$ at $p$ and $\mathbf{x}_2$ at $\varphi(p)$
  and let $\tilde\varphi=\mathbf{x}_2\circ\varphi\circ\mathbf{x}_1^{-1}$ be the
  coordinate representative, a map between open subsets of $\mathbb{R}^n$. Because
  $\varphi$ is differentiable, $\tilde\varphi$ is $C^\infty$ on a neighbourhood of
  $q=\mathbf{x}_1(p)$, and its Fréchet derivative $d\tilde\varphi_q$ is the coordinate
  expression of $d\varphi_p$, hence a linear isomorphism. The inverse function theorem
  in $\mathbb{R}^n$ then makes $\tilde\varphi$ a local diffeomorphism at $q$: the
  analytic core is that a $C^\infty$ map whose derivative at a point is invertible has
  a $C^\infty$ local inverse, the derivative being invertible on a whole neighbourhood
  and the inverse's smoothness propagating pointwise. Finally $\varphi=\mathbf{x}_2^{-1}
  \circ\tilde\varphi\circ\mathbf{x}_1$ near $p$ is a composition of local
  diffeomorphisms (the two charts and $\tilde\varphi$), hence a local diffeomorphism
  at $p$.
\end{proof}

\section{Immersions, embeddings, and submanifolds}

\begin{definition}[immersion; embedding; submanifold]
  \label{def:dc-ch0-3-1}
  \lean{Riemannian.DCIsImmersion, Riemannian.DCIsEmbedding}
  \uses{def:dc-ch0-2-5, def:dc-ch0-2-8}
  \leanok
  \dcref{ch0:3.1}
  % submanifold = image of an embedded inclusion (an instance of DCIsEmbedding);
  % codimension $n-m$ is a derived numeric notion, not separate Lean data.
  Let $M^m,N^n$ be differentiable manifolds. A differentiable $\varphi:M\to N$ is an
  \emph{immersion} if $d\varphi_p:T_pM\to T_{\varphi(p)}N$ is injective for all $p\in M$.
  If in addition $\varphi$ is a homeomorphism onto $\varphi(M)\subset N$ (subspace
  topology), $\varphi$ is an \emph{embedding}. If $M\subset N$ and the inclusion
  $i:M\subset N$ is an embedding, $M$ is a \emph{submanifold} of $N$. If $\varphi:M^m\to
  N^n$ is an immersion then $m\le n$; the difference $n-m$ is the \emph{codimension}.
\end{definition}

\begin{example}[a regular surface in $\mathbb{R}^3$ is a submanifold]
  \label{ex:dc-ch0-3-6}
  \notready
  \dcref{ch0:3.6}
  \notready
  A regular surface $S\subset\mathbb{R}^3$ has the differentiable structure given by
  its parametrizations $\mathbf{x}_\alpha:U_\alpha\to S$; these are embeddings of
  $U_\alpha$ into $S$, by conditions (a),(b) of the definition of regular surface.
  The inclusion $i:S\subset\mathbb{R}^3$ is differentiable (for $p\in S$ take
  $\mathbf{x}$ at $p$ and $j:V\subset\mathbb{R}^3\to V$ the identity, so
  $j^{-1}\circ i\circ\mathbf{x}=\mathbf{x}$ is differentiable); from (b) $i$ is an
  immersion and from (a) $i$ is a homeomorphism onto its image. Hence $S$ is a
  submanifold of $\mathbb{R}^3$.
\end{example}

\begin{proposition}[every immersion is locally an embedding]
  \label{prop:dc-ch0-3-7}
  \lean{Riemannian.DCIsImmersion.exists_isEmbedding_restrict}
  \uses{def:dc-ch0-3-1, lem:dc-ch0-3-7-embedding, lem:dc-ch0-3-7-normalform}
  \leanok
  \dcref{ch0:3.7}
  Let $\varphi:M_1^n\to M_2^m$, $n\le m$, be an immersion. For every $p\in M_1$
  there exists a neighborhood $V\subset M_1$ of $p$ such that the restriction
  $\varphi\,|\,V\to M_2$ is an embedding.
\end{proposition}
\begin{proof}
  \sketch
  \leanok
  A consequence of the inverse function theorem. Take coordinates
  $\mathbf{x}_1:U_1\to M_1$ at $p$, $\mathbf{x}_2:U_2\to M_2$ at $\varphi(p)$, and
  write $\tilde\varphi=\mathbf{x}_2^{-1}\circ\varphi\circ\mathbf{x}_1=(y_1(x),\dots,y_m(x))$.
  With $q=\mathbf{x}_1^{-1}(p)$, renumber so that
  $\frac{\partial(y_1,\dots,y_n)}{\partial(x_1,\dots,x_n)}(q)\neq0$. Extend to
  $\phi:U_1\times\mathbb{R}^{m-n=k}\to\mathbb{R}^m$ by
  
  $$
  \phi(x_1,\dots,x_n,t_1,\dots,t_k)=(y_1(x),\dots,y_n(x),y_{n+1}(x)+t_1,\dots,y_{n+k}(x)+t_k),
  $$
  
  which restricts to $\tilde\varphi$ on $U_1\times\{0\}$ and has
  $\det(d\phi_q)=\frac{\partial(y_1,\dots,y_n)}{\partial(x_1,\dots,x_n)}(q)\neq0$. By
  the inverse function theorem $\phi\,|\,W_1$ is a diffeomorphism onto $W_2$; with
  $\tilde V=W_1\cap U_1$, $V=\mathbf{x}_1(\tilde V)$, the restriction
  $\mathbf{x}_2\circ\tilde\varphi\circ\mathbf{x}_1^{-1}:V\to\varphi(V)$ is a
  diffeomorphism, hence an embedding. This is the normal-form argument of
  \cref{lem:dc-ch0-3-7-normalform}.

\end{proof}

\begin{lemma}[immersion normal form yields a local embedding]
  \label{lem:dc-ch0-3-7-embedding}
  \lean{Riemannian.isEmbedding_restrict_of_isImmersionAt}
  \uses{def:dc-ch0-3-1}
  \leanok
  \dcref{ch0:3.7}
  Let $\varphi:M_1\to M_2$ and $p\in M_1$. Suppose that in suitable coordinate charts
  around $p$ and $\varphi(p)$ the map $\varphi$ has the standard immersion normal form
  $u\mapsto(u,0)$ (a linear inclusion of the model space of $M_1$ into that of $M_2$).
  Then there is an open neighborhood $U\ni p$ such that the restriction $\varphi\,|\,U$ is
  a topological embedding. Hence every immersion is \emph{locally} an embedding.
\end{lemma}
\begin{proof}
  \sketch
  \leanok
  In the chart around $p$, $\varphi$ is the composite of three maps: the chart, which is a
  homeomorphism of a neighborhood of $p$ onto an open subset of the model space; the linear
  inclusion $u\mapsto(u,0)$, which is a topological embedding; and the inverse of the chart
  around $\varphi(p)$, a homeomorphism onto its image. The subspace topology on $U$ is
  therefore induced by $\varphi$ (each factor induces its topology, and inducing maps
  compose), and $\varphi\,|\,U$ is injective because the chart, the linear inclusion, and
  the codomain chart are each injective. An injective inducing map is an embedding.
\end{proof}

\begin{lemma}[local immersion normal form in normed spaces]
  \label{lem:dc-ch0-3-7-normalform}
  \lean{Riemannian.isEmbedding_restrict_of_hasFDerivAt_injective}
  \leanok
  \dcref{ch0:3.7}
  Let $g:E\to F$ be a $C^\infty$ map between finite-dimensional normed spaces, defined on an
  open set $s\ni q$, whose Fréchet derivative $L=dg_q$ is injective. Then $g$ restricts to a
  topological embedding on a neighborhood of $q$.
\end{lemma}
\begin{proof}
  \sketch
  \leanok
  Choose a complement $G$ of the range of $L$ in $F$ (finite-dimensional, hence closed and
  complemented), so that $(x,t)\mapsto Lx+t$ is a linear isomorphism $E\times G\to F$. Extend
  $g$ to $\mathrm{Ext}(x,t)=g(x)+t$ on $s\times G$; its derivative at $(q,0)$ is exactly that
  isomorphism, so by the inverse function theorem $\mathrm{Ext}$ is a homeomorphism from an
  open neighborhood of $(q,0)$ onto an open subset of $F$. The slice inclusion
  $u\mapsto(u,0)$ is a topological embedding, and $g$ agrees with $\mathrm{Ext}\circ(u\mapsto(u,0))$
  near $q$; being a composite of an embedding with a homeomorphism onto its image, $g$ is a
  topological embedding on a neighborhood of $q$.
\end{proof}

\section{Other examples of manifolds. Orientation}

\begin{example}[the tangent bundle]
  \label{ex:dc-ch0-4-1}
  \lean{Riemannian.DCTangentBundle}
  \uses{def:dc-ch0-2-1}
  \leanok
  \dcref{ch0:4.1}
  Let $M^n$ be a differentiable manifold and $TM=\{(p,v); p\in M, v\in T_pM\}$. With
  a maximal structure $\{(U_\alpha,\mathbf{x}_\alpha)\}$ on $M$, coordinates
  $(x_1^\alpha,\dots,x_n^\alpha)$ on $U_\alpha$ and bases
  $\{\frac{\partial}{\partial x_i^\alpha}\}$, define
  $\mathbf{y}_\alpha:U_\alpha\times\mathbb{R}^n\to TM$ by
  $$
  \mathbf{y}_\alpha(x_1^\alpha,\dots,x_n^\alpha,u_1,\dots,u_n)=\Big(\mathbf{x}_\alpha(x_1^\alpha,\dots,x_n^\alpha),\sum_{i=1}^n u_i\frac{\partial}{\partial x_i^\alpha}\Big).
  $$
  Then $\bigcup_\alpha\mathbf{y}_\alpha(U_\alpha\times\mathbb{R}^n)=TM$ verifies (1)
  of \cref{def:dc-ch0-2-1}, and since $\mathbf{x}_\beta^{-1}\circ\mathbf{x}_\alpha$ is
  differentiable so is $d(\mathbf{x}_\beta^{-1}\circ\mathbf{x}_\alpha)$, giving
  $$
  \mathbf{y}_\beta^{-1}\circ\mathbf{y}_\alpha(q_\alpha,v_\alpha)=\big((\mathbf{x}_\beta^{-1}\circ\mathbf{x}_\alpha)(q_\alpha),\,d(\mathbf{x}_\beta^{-1}\circ\mathbf{x}_\alpha)(v_\alpha)\big)
  $$
  differentiable, verifying (2). Thus $\{(U_\alpha\times\mathbb{R}^n,\mathbf{y}_\alpha)\}$
  is a differentiable structure of dimension $2n$ on $TM$, the \emph{tangent bundle}.
\end{example}

\begin{lemma}[the tangent bundle is a smooth manifold]
  \label{lem:dc-ch0-4-1-manifold}
  \lean{Riemannian.DCTangentBundle.isManifold}
  \uses{ex:dc-ch0-4-1}
  \leanok
  \dcref{ch0:4.1}
  The tangent bundle $TM$ of a smooth $n$-manifold $M$ carries a smooth manifold
  structure of dimension $2n$, modelled on $H\times E$.
\end{lemma}
\begin{proof}
  \sketch
  \leanok
  The charts $\mathbf{y}_\alpha$ of \cref{ex:dc-ch0-4-1} form a differentiable
  atlas: the transition maps
  $\mathbf{y}_\beta^{-1}\circ\mathbf{y}_\alpha(q,v)=\big((\mathbf{x}_\beta^{-1}\circ\mathbf{x}_\alpha)(q),\,d(\mathbf{x}_\beta^{-1}\circ\mathbf{x}_\alpha)(v)\big)$
  are smooth, being built from the smooth change of coordinates on $M$ and its
  differential. Hence $TM$ is a smooth manifold modelled on $H\times E$.
\end{proof}

\begin{example}[regular surfaces in $\mathbb{R}^n$]
  \label{ex:dc-ch0-4-2}
  \uses{ex:dc-ch0-3-6}
  \notready
  \dcref{ch0:4.2}
  \notready
  $M^k\subset\mathbb{R}^n$, $k\le n$, is a \emph{regular surface of dimension $k$} if for
  every $p\in M^k$ there is a neighborhood $V$ of $p$ in $\mathbb{R}^n$ and
  $\mathbf{x}:U\subset\mathbb{R}^k\to M\cap V$ with (a) $\mathbf{x}$ a
  differentiable homeomorphism and (b) $(d\mathbf{x})_q:\mathbb{R}^k\to\mathbb{R}^n$
  injective for all $q\in U$. If
  $\mathbf{x}:U\to M^k$, $\mathbf{y}:V\to M^k$ are two parametrizations with
  $\mathbf{x}(U)\cap\mathbf{y}(V)=W\neq\phi$ then
  $h=\mathbf{x}^{-1}\circ\mathbf{y}:\mathbf{y}^{-1}(W)\to\mathbf{x}^{-1}(W)$ is a
  diffeomorphism. Consequently, $M^k$ is a differentiable manifold of dimension
  $k$, and the inclusion $i:M^k\subset\mathbb{R}^n$ is an embedding.
\end{example}
\begin{proof}
  \sketch
  $h$ is a homeomorphism (composition of homeomorphisms). For
  $r\in\mathbf{y}^{-1}(W)$, $q=h(r)$, write
  $\mathbf{x}(u_1,\dots,u_k)=(v_1,\dots,v_n)$; by (b) suppose
  $\frac{\partial(v_1,\dots,v_k)}{\partial(u_1,\dots,u_k)}(q)\neq0$. Extend
  $\mathbf{x}$ to $F:U\times\mathbb{R}^{n-k}\to\mathbb{R}^n$,
  $F(u,t)=(v_1(u),\dots,v_k(u),v_{k+1}(u)+t_{k+1},\dots,v_n(u)+t_n)$, with
  $\det(dF_q)=\frac{\partial(v_1,\dots,v_k)}{\partial(u_1,\dots,u_k)}(q)\neq0$. By
  the inverse function theorem $F^{-1}$ exists and is differentiable on a
  neighborhood $Q$ of $\mathbf{x}(q)$; choosing $R$ with $\mathbf{y}(R)\subset Q$,
  $h\,|\,R=F^{-1}\circ\mathbf{y}\,|\,R$ is differentiable, hence $h$ is, and
  similarly $h^{-1}$. The final assertion follows from
  \cref{ex:dc-ch0-3-6}.
\end{proof}

\begin{example}[inverse image of a regular value]
  \label{ex:dc-ch0-4-3}
  \uses{ex:dc-ch0-4-2}
  \notready
  \dcref{ch0:4.3}
  \notready
  Let $F:U\subset\mathbb{R}^n\to\mathbb{R}^m$ be differentiable on an open set $U$.
  A point $p\in U$ is a \emph{critical point} if $dF_p:\mathbb{R}^n\to\mathbb{R}^m$ is
  not surjective; $F(p)$ is then a \emph{critical value}. A point $a\in\mathbb{R}^m$ not
  a critical value is a \emph{regular value} (any $a\notin F(U)$ is trivially regular; if
  a regular value of $F$ exists in $\mathbb{R}^m$ then $n\ge m$). If $a\in F(U)$ is
  a regular value, then $F^{-1}(a)\subset\mathbb{R}^n$ is a regular surface of
  dimension $n-m=k$, hence (\cref{ex:dc-ch0-4-2}) a submanifold of $\mathbb{R}^n$.
\end{example}
\begin{proof}
  \sketch
  For $p\in F^{-1}(a)$, writing $q=(y_1,\dots,y_m,x_1,\dots,x_k)$ and
  $F(q)=(f_1(q),\dots,f_m(q))$, since $dF_p$ is surjective suppose
  $\frac{\partial(f_1,\dots,f_m)}{\partial(y_1,\dots,y_m)}(p)\neq0$. Define
  $\varphi:U\to\mathbb{R}^{m+k}$ by
  $\varphi(y,x)=(f_1(q),\dots,f_m(q),x_1,\dots,x_k)$, with
  $\det(d\varphi)_p=\frac{\partial(f_1,\dots,f_m)}{\partial(y_1,\dots,y_m)}(p)\neq0$.
  By the inverse function theorem $\varphi$ is a diffeomorphism of a neighborhood
  $Q$ of $p$ onto $W$; with a cube $K^{m+k}$ centered at $\varphi(p)$ and
  $V=\varphi^{-1}(K^{m+k})\cap Q$, $\varphi$ maps $V$ diffeomorphically onto
  $K^m\times K^k$, and $\mathbf{x}:K^k\to V$,
  $\mathbf{x}(x_1,\dots,x_k)=\varphi^{-1}(a_1,\dots,a_m,x_1,\dots,x_k)$, satisfies
  (a),(b) of \cref{ex:dc-ch0-4-2}. Hence $F^{-1}(a)$ is a regular surface.
\end{proof}

\begin{definition}[orientable manifold; orientation]
  \label{def:dc-ch0-4-4}
  \lean{Riemannian.DCIsOrientable}
  \leanok
  \dcref{ch0:4.4}
  $M$ is \emph{orientable} if it admits a differentiable structure
  $\{(U_\alpha,\mathbf{x}_\alpha)\}$ such that:
  (i) for every pair $\alpha,\beta$ with
  $\mathbf{x}_\alpha(U_\alpha)\cap\mathbf{x}_\beta(U_\beta)=W\neq\phi$, the
  differential of the change of coordinates
  $\mathbf{x}_\beta^{-1}\circ\mathbf{x}_\alpha$ has positive determinant.
  Otherwise $M$ is \emph{non-orientable}. A choice of structure satisfying (i) is an
  \emph{orientation}; $M$ is then \emph{oriented}. Two structures satisfying (i) \emph{determine
  the same orientation} if their union again satisfies (i). If $M$ is orientable and
  connected there exist exactly two orientations. For a diffeomorphism
  $\varphi:M_1\to M_2$, $M_1$ is orientable iff $M_2$ is; if both are connected and
  oriented, $\varphi$ \emph{preserves} or \emph{reverses} the orientation according as the
  induced orientation coincides with the given one or not.
\end{definition}

\begin{example}[two-neighborhood criterion]
  \label{ex:dc-ch0-4-5}
  \notready
  \dcref{ch0:4.5}
  \notready
  If $M$ can be covered by two coordinate neighborhoods $V_1,V_2$ with
  $V_1\cap V_2$ connected, then $M$ is orientable: the determinant of the change of
  coordinates has constant sign on the connected $V_1\cap V_2$, and if negative one
  changes the sign of one coordinate to make it positive throughout.
\end{example}

\begin{example}[another description of projective space]
  \label{ex:dc-ch0-4-7}
  \uses{ex:dc-ch0-2-4}
  \notready
  \dcref{ch0:4.7}
  \notready
  $P^n(\mathbb{R})$ is the quotient of the unit sphere $S^n=\{p\in\mathbb{R}^{n+1};
  |p|=1\}$ by $p\sim A(p)=-p$. Introducing on $S^n\subset\mathbb{R}^{n+1}$ the
  regular-surface structure with hemisphere parametrizations
  $\mathbf{x}_i^{\pm}:U_i\to S^n$,
  $U_i=\{x_i=0,\,x_1^2+\dots+x_{i-1}^2+x_{i+1}^2+\dots+x_{n+1}^2<1\}$,
  $\mathbf{x}_i^{\pm}(x_1,\dots,\hat x_i,\dots,x_{n+1})=(x_1,\dots,\pm D_i,\dots,x_{n+1})$,
  $D_i=\sqrt{1-(x_1^2+\dots+x_{i-1}^2+x_{i+1}^2+\dots+x_{n+1}^2)}$. With the
  canonical projection $\pi:S^n\to P^n(\mathbb{R})$, $\pi(p)=\{p,-p\}$, define
  $\mathbf{y}_i=\pi\circ\mathbf{x}_i^+:U_i\to P^n(\mathbb{R})$. Then
  $\mathbf{y}_i^{-1}\circ\mathbf{y}_j=(\mathbf{x}_i^+)^{-1}\circ\mathbf{x}_j^+$ is
  differentiable, so $\{(U_i,\mathbf{y}_i)\}$ is a differentiable structure for
  $P^n(\mathbb{R})$, giving the same maximal structure as \cref{ex:dc-ch0-2-4}.
  The projective space is orientable exactly when $n$ is odd.
\end{example}

\begin{example}[discontinuous action of a group]
  \label{ex:dc-ch0-4-8}
  \uses{ex:dc-ch0-4-7}
  \notready
  \dcref{ch0:4.8}
  \notready
  A group $G$ \emph{acts on} a differentiable manifold $M$ if there is
  $\varphi:G\times M\to M$ such that (i) for each $g\in G$, $\varphi_g(p)=\varphi(g,p)$
  is a diffeomorphism with $\varphi_e=\mathrm{id}$, and (ii)
  $\varphi_{g_1 g_2}=\varphi_{g_1}\circ\varphi_{g_2}$ (writing $gp$ for $\varphi(g,p)$,
  this is $(g_1 g_2)p=g_1(g_2 p)$). The action is \emph{properly discontinuous} if every
  $p\in M$ has a neighborhood $U$ with $U\cap g(U)=\phi$ for all $g\neq e$. The
  action determines $\sim$ on $M$ ($p_1\sim p_2$ iff $p_2=gp_1$), with quotient
  $M/G$ and \emph{projection} $\pi:M\to M/G$, $\pi(p)=Gp$.
  For a properly discontinuous action, $M/G$ has a differentiable structure with
  respect to which $\pi:M\to M/G$ is a local diffeomorphism.
\end{example}
\begin{proof}
  \sketch
  For $p\in M$ choose $\mathbf{x}:V\to M$ at $p$ with $\mathbf{x}(V)\subset
  U$, $U\cap g(U)=\phi$ for $g\neq e$; then $\pi\,|\,U$ is injective, so
  $\mathbf{y}=\pi\circ\mathbf{x}:V\to M/G$ is injective and $\{(V,\mathbf{y})\}$
  covers $M/G$. Given $\mathbf{y}_1=\pi\circ\mathbf{x}_1$,
  $\mathbf{y}_2=\pi\circ\mathbf{x}_2$ with overlapping images, with $\pi_i=\pi\,|\,\mathbf{x}_i(V_i)$,
  on a suitable $W$ the change is
  $\mathbf{y}_1^{-1}\circ\mathbf{y}_2\,|\,W=\mathbf{x}_1^{-1}\circ\pi_1^{-1}\circ\pi_2\circ\mathbf{x}_2$;
  it suffices that $\pi_1^{-1}\circ\pi_2$ be differentiable at
  $p_2=\pi_2^{-1}(q)$. With $p_1=\pi_1^{-1}\circ\pi_2(p_2)$ there is $g\in G$ with
  $gp_2=p_1$, and $\pi_1^{-1}\circ\pi_2\,|\,\mathbf{x}_2(W)$ coincides with the
  diffeomorphism $\varphi_g\,|\,\mathbf{x}_2(W)$. Hence $\pi$ is a local
  diffeomorphism. The quotient in \cref{ex:dc-ch0-4-7} is obtained from
  $M=S^n$ and $G=\{A,I=A^2\}$.
\end{proof}

\section{Vector fields; brackets. Topology of manifolds}

\begin{definition}[vector field]
  \label{def:dc-ch0-5-1}
  \lean{Riemannian.VectorFieldSection, Riemannian.TangentSmoothAt, Riemannian.TangentSmoothAt.iff_coord, Riemannian.SmoothVectorField.smoothAt}
  \uses{ex:dc-ch0-4-1}
  \leanok
  \dcref{ch0:5.1}
  A \emph{vector field $X$ on $M$} associates to each $p\in M$ a vector $X(p)\in T_pM$;
  as a mapping, $X:M\to TM$. $X$ is \emph{differentiable} if $X:M\to TM$
  is. In a parametrization $\mathbf{x}:U\subset\mathbb{R}^n\to M$,
  $$
  X(p)=\sum_{i=1}^n a_i(p)\frac{\partial}{\partial x_i},\qquad\text{(4)}
  $$
  $a_i:U\to\mathbb{R}$; $X$ is differentiable iff the $a_i$ are differentiable for
  some (hence any) parametrization. Let $\mathcal{D}$ be the differentiable
  functions on $M$ and $\mathcal{F}$ the functions on $M$; thinking of
  $X:\mathcal{D}\to\mathcal{F}$,
  $$
  (Xf)(p)=\sum_i a_i(p)\frac{\partial f}{\partial x_i}(p),\qquad\text{(5)}
  $$
  $Xf$ is independent of $\mathbf{x}$, and $X$ is differentiable iff
  $Xf\in\mathcal{D}$ for all $f\in\mathcal{D}$. If $\varphi:M\to M$ is a
  diffeomorphism, $v\in T_pM$, and $f$ differentiable near $\varphi(p)$, then
  $(d\varphi(v)f)\varphi(p)=v(f\circ\varphi)(p)$.
\end{definition}

\begin{lemma}[existence and uniqueness of the bracket field]
  \label{lem:dc-ch0-5-2}
  \lean{Riemannian.exists_unique_bracketField, Riemannian.bracketField_dir}
  \leanok
  \dcref{ch0:5.2}
  Let $X,Y$ be differentiable vector fields on $M$. Then there exists a unique
  differentiable vector field $Z$ such that $Zf=(XY-YX)f$ for all
  $f\in\mathcal{D}$. This field is the \emph{bracket} $[X,Y]=XY-YX$.
  Existence, uniqueness, and smoothness are isolated in
  \cref{lem:dc-ch0-5-2-commutator,lem:dc-ch0-5-2-unique,lem:dc-ch0-5-2-smooth};
  the structural identities are treated in \cref{lem:dc-ch0-5-3-antisymm,lem:dc-ch0-5-3-jacobi}.
\end{lemma}
\begin{proof}
  \sketch
  \leanok
  \uses{lem:dc-ch0-5-2-unique,lem:dc-ch0-5-2-commutator,lem:dc-ch0-5-2-smooth}
  In a parametrization $\mathbf{x}$ at $p$ with
  $X=\sum_i a_i\frac{\partial}{\partial x_i}$,
  $Y=\sum_j b_j\frac{\partial}{\partial x_j}$,
  
  $$
  XYf=\sum_{i,j}a_i\frac{\partial b_j}{\partial x_i}\frac{\partial f}{\partial x_j}+\sum_{i,j}a_i b_j\frac{\partial^2 f}{\partial x_i\partial x_j},
  $$
  
  and symmetrically for $YXf$; the second-order terms cancel, giving
  
  $$
  Zf=XYf-YXf=\sum_{i,j}\Big(a_i\frac{\partial b_j}{\partial x_i}-b_i\frac{\partial a_j}{\partial x_i}\Big)\frac{\partial f}{\partial x_j},
  $$
  
  proving uniqueness. For existence, define $Z_\alpha$ on each
  $\mathbf{x}_\alpha(U_\alpha)$ by this expression; by uniqueness $Z_\alpha=Z_\beta$
  on overlaps, so $Z$ is defined over all of $M$. The coordinate formula shows
  that it is differentiable.
\end{proof}

\begin{lemma}[existence characterization: the bracket acts as the commutator of derivations]
  \label{lem:dc-ch0-5-2-commutator}
  \lean{Riemannian.mfderiv_mlieBracket_eq_commutator}
  \uses{lem:dc-ch0-5-2}
  \leanok
  \dcref{ch0:5.2}
  Let $X,Y$ be smooth vector fields on a boundaryless manifold $M$ and $f$ a smooth
  real function. Then the bracket field $[X,Y]$ acts on $f$ as the
  commutator of the derivations $X$ and $Y$:
  $$
  [X,Y]f = X(Yf)-Y(Xf), \qquad\text{i.e.}\qquad
  \mathrm{d}f_p([X,Y]_p) = \mathrm{d}(\mathrm{d}f\,Y)_p(X_p)-\mathrm{d}(\mathrm{d}f\,X)_p(Y_p).
  $$
  Thus $[X,Y]$ realizes the derivation $(XY-YX)$.
\end{lemma}
\begin{proof}
  \sketch
  \leanok
  In a chart around $p$, the boundaryless hypothesis identifies the model set with the
  model space. Pull back $X,Y$ to this chart and apply the Euclidean commutator identity
  $\mathrm{d}f([V,W])=\mathrm{d}(\mathrm{d}f\,W)(V)-\mathrm{d}(\mathrm{d}f\,V)(W)$;
  the chain rule transports the result back to $M$.
\end{proof}

\begin{lemma}[uniqueness: a vector field is determined by its action on functions]
  \label{lem:dc-ch0-5-2-unique}
  \lean{Riemannian.tangentSection_eq_of_forall_mfderiv}
  \uses{lem:dc-ch0-5-2}
  \leanok
  \dcref{ch0:5.2}
  A tangent vector at $p$, and hence a vector field on $M$, is determined by its
  action on differentiable real functions: if $Z_1,Z_2$ are sections of $TM$ with
  $Z_1f=Z_2f$ for every $f\in\mathcal{D}$, then $Z_1=Z_2$; in particular at most one
  field acts as $(XY-YX)$.
\end{lemma}
\begin{proof}
  \sketch
  \leanok
  It suffices to prove the pointwise statement: if $v\in T_pM$ satisfies
  $\mathrm{d}f_p(v)=0$ for every $f$ differentiable at $p$, then $v=0$
  (apply it to $Z_1(p)-Z_2(p)$ using linearity of $\mathrm{d}f_p$). For a
  continuous linear functional $L$ on the model space $E$, the coordinate function
  $f_L=L\circ\mathbf{x}$
  is differentiable at $p$, and since $\mathrm{d}\mathbf{x}_p=\mathrm{id}$ on
  $T_pM$ one has $\mathrm{d}(f_L)_p(v)=L(v)$. Thus $L(v)=0$ for all such $L$, and
  continuous functionals separate points of a normed space, so $v=0$.
\end{proof}

\begin{lemma}[smoothness of the bracket field]
  \label{lem:dc-ch0-5-2-smooth}
  \lean{Riemannian.DCLieBracket_smoothAt}
  \uses{lem:dc-ch0-5-2}
  \leanok
  \dcref{ch0:5.2}
  If $X$ and $Y$ are smooth vector fields, then their bracket $[X,Y]$ is a smooth
  vector field.
\end{lemma}
\begin{proof}
  \sketch
  \leanok
  In each parametrization, write
  $X=\sum_i a_i\frac{\partial}{\partial x_i}$ and
  $Y=\sum_j b_j\frac{\partial}{\partial x_j}$. By
  \cref{lem:dc-ch0-5-2}, the bracket has coefficients
  $\sum_i(a_i\partial_i b_j-b_i\partial_i a_j)$. These coefficients are smooth
  because the coefficients of $X$ and $Y$ are smooth, so the local expressions
  define a smooth vector field on each chart. Since the bracket was constructed
  uniquely, the local smooth fields agree on overlaps and glue to a smooth global
  vector field.
\end{proof}

\begin{lemma}[antisymmetry of the bracket]
  \label{lem:dc-ch0-5-3-antisymm}
  \lean{Riemannian.DCLieBracket_antisymm}
  \uses{lem:dc-ch0-5-2}
  \leanok
  \dcref{ch0:5.3}
  For differentiable vector fields $X$ and $Y$ one has $[X,Y]=-[Y,X]$.
\end{lemma}
\begin{proof}
  \sketch
  \leanok
  By the defining identity of \cref{lem:dc-ch0-5-2},
  $[X,Y]f=XYf-YXf$ and $[Y,X]f=YXf-XYf$ for every differentiable function $f$.
  Hence $[X,Y]f=-[Y,X]f$ for all such $f$, and uniqueness in
  \cref{lem:dc-ch0-5-2} gives $[X,Y]=-[Y,X]$.
\end{proof}

\begin{lemma}[Jacobi identity in Leibniz form]
  \label{lem:dc-ch0-5-3-jacobi}
  \lean{Riemannian.DCLieBracket_jacobi}
  \uses{lem:dc-ch0-5-2}
  \leanok
  \dcref{ch0:5.3}
  For differentiable vector fields $X,Y,Z$ one has
  $[X,[Y,Z]]=[[X,Y],Z]+[Y,[X,Z]]$.
\end{lemma}
\begin{proof}
  \sketch
  \leanok
  Evaluate both sides on a differentiable function $f$. Expanding each bracket
  as a commutator of derivations gives
  $XYZf-XZYf-YZXf+ZYXf$ on the left and the same four terms on the right after
  cancellation of the repeated middle terms. By uniqueness in
  \cref{lem:dc-ch0-5-2}, the two vector fields are equal.
\end{proof}

\begin{lemma}[Jacobi identity in cyclic form]
  \label{lem:dc-ch0-5-3-jacobi-cyclic}
  \lean{Riemannian.DCLieBracket_jacobi_cyclic}
  \uses{lem:dc-ch0-5-3-jacobi, lem:dc-ch0-5-3-antisymm}
  \leanok
  \dcref{ch0:5.3}
  For differentiable vector fields $X,Y,Z$ one has the cyclic Jacobi identity
  $[[X,Y],Z]+[[Y,Z],X]+[[Z,X],Y]=0$.
\end{lemma}
\begin{proof}
  \sketch
  \leanok
  Anticommutativity (\cref{lem:dc-ch0-5-3-antisymm}) rewrites each outer bracket as
  $[[X,Y],Z]=-[Z,[X,Y]]$, so the cyclic sum equals
  $-\big([Z,[X,Y]]+[X,[Y,Z]]+[Y,[Z,X]]\big)$. The Leibniz-form identity
  $[X,[Y,Z]]=[[X,Y],Z]+[Y,[X,Z]]$ of \cref{lem:dc-ch0-5-3-jacobi}, together with
  $[X,Z]=-[Z,X]$, makes the parenthesized derivation sum vanish, giving $0$.
\end{proof}

\begin{lemma}[linearity of the bracket]
  \label{lem:dc-ch0-5-3-bilinear}
  \lean{Riemannian.DCLieBracket_linear_left}
  \uses{lem:dc-ch0-5-2}
  \leanok
  \dcref{ch0:5.3}
  For smooth vector fields $X,Y,Z$ and scalars $a,b\in\mathbb{R}$ one has
  $[aX+bY,Z]=a[X,Z]+b[Y,Z]$.
\end{lemma}
\begin{proof}
  \sketch
  \leanok
  Apply additivity of the bracket in its first slot to the sum $aX+bY$, then pull
  out the constants $a$ and $b$ by homogeneity in the first slot. Both facts follow
  from the local commutator expression of \cref{lem:dc-ch0-5-2}.
\end{proof}

\begin{lemma}[Leibniz rule for the bracket]
  \label{lem:dc-ch0-5-3-leibniz}
  \lean{Riemannian.DCLieBracket_leibniz}
  \uses{lem:dc-ch0-5-2}
  \leanok
  \dcref{ch0:5.3}
  For smooth vector fields $X,Y$ and smooth functions $f,g$ one has
  $[fX,gY]=fg[X,Y]+fX(g)Y-gY(f)X$, where $X(g)=dg_p(X_p)$ denotes the derivative
  of $g$ along $X$.
\end{lemma}
\begin{proof}
  \sketch
  \leanok
  Expand the bracket by the first-slot Leibniz identity
  $[fX,W]=f[X,W]-W(f)X$ with $W=gY$, then the second-slot Leibniz identity
  $[X,gY]=g[X,Y]+X(g)Y$; combining and using linearity of the differential
  $dg_p$ gives $[fX,gY]=fg[X,Y]+fX(g)Y-gY(f)X$.
\end{proof}

\begin{proposition}[properties of the bracket]
  \label{prop:dc-ch0-5-3}
  \lean{Riemannian.DCLieBracket_properties}
  \uses{lem:dc-ch0-5-3-antisymm, lem:dc-ch0-5-3-jacobi-cyclic, lem:dc-ch0-5-3-bilinear, lem:dc-ch0-5-3-leibniz}
  \leanok
  \dcref{ch0:5.3}
  If $X,Y,Z$ are differentiable vector fields on $M$, $a,b\in\mathbb{R}$, and $f,g$
  differentiable functions, then:
  (a) $[X,Y]=-[Y,X]$ \emph{(anticommutativity)};
  (b) $[aX+bY,Z]=a[X,Z]+b[Y,Z]$ \emph{(linearity)};
  (c) $[[X,Y],Z]+[[Y,Z],X]+[[Z,X],Y]=0$ \emph{(Jacobi identity)};
  (d) $[fX,gY]=fg[X,Y]+fX(g)Y-gY(f)X$.
\end{proposition}
\begin{proof}
  \sketch
  (a),(b) are immediate. For (c), expand
  $[[X,Y],Z]=XYZ-YXZ-ZXY+ZYX$ and
  $[X,[Y,Z]]+[Y,[Z,X]]=XYZ-XZY-YZX+ZYX+YZX-YXZ-ZXY+XZY$; the second members
  agree, and (c) follows using (a). For (d),
  $[fX,gY]=fX(gY)-gY(fX)=fgXY+fX(g)Y-gfYX-gY(f)X=fg[X,Y]+fX(g)Y-gY(f)X$.
\end{proof}

For a differentiable vector field $X$ and $p\in M$, there exist a neighborhood
$U$ of $p$, an interval $(-\delta,\delta)$, and a differentiable map
$\varphi:(-\delta,\delta)\times U\to M$ such that $t\mapsto\varphi(t,q)$ is the
unique curve with
$\frac{\partial\varphi}{\partial t}=X(\varphi(t,q))$ and $\varphi(0,q)=q$. A
curve $\alpha$ with $\alpha'(t)=X(\alpha(t))$, $\alpha(0)=q$, is a
\emph{trajectory} of $X$, and $\varphi_t(q)=\varphi(t,q)$ is the \emph{local
flow} of $X$.

\begin{proposition}[bracket as Lie derivative along the flow]
  \label{prop:dc-ch0-5-4}
  \lean{Riemannian.DCLieBracket_eq_flow_lieDerivative}
  \uses{lem:dc-ch0-5-5, lem:dc-ch0-5-2}
  \leanok
  \dcref{ch0:5.4}
  Let $X,Y$ be differentiable vector fields on $M$, $p\in M$, and $\varphi_t$ the
  local flow of $X$ in a neighborhood $U$ of $p$. Then

  $$
  [X,Y](p)=\lim_{t\to0}\frac1t\,[Y-d\varphi_t Y](\varphi_t(p)).
  $$
\end{proposition}
\begin{proof}
  \sketch
  \leanok
  For $f$ differentiable near $p$, put $h(t,q)=f(\varphi_t(q))-f(q)$; by
  \cref{lem:dc-ch0-5-5} there is differentiable $g(t,q)$ with
  $f\circ\varphi_t(q)=f(q)+tg(t,q)$ and $g(0,q)=Xf(q)$. Then
  $((d\varphi_t Y)f)(\varphi_t(p))=(Y(f\circ\varphi_t))(p)=Yf(p)+t(Yg(t,p))$, so

  $$
  \lim_{t\to0}\frac1t[Y-d\varphi_t Y]f(\varphi_t p)=\lim_{t\to0}\frac{(Yf)(\varphi_t p)-Yf(p)}{t}-(Yg(0,p))=(X(Yf))(p)-(Y(Xf))(p)=([X,Y]f)(p).
  $$
\end{proof}

\begin{lemma}[Hadamard-type lemma from calculus]
  \label{lem:dc-ch0-5-5}
  \lean{Riemannian.DCHadamardFactor, Riemannian.DCHadamardFactorDiff}
  \leanok
  \dcref{ch0:5.5}
  Let $h:(-\delta,\delta)\times U\to\mathbb{R}$ be differentiable with $h(0,q)=0$
  for all $q\in U$. Then there exists a differentiable
  $g:(-\delta,\delta)\times U\to\mathbb{R}$ with $h(t,q)=t\,g(t,q)$; in particular
  $g(0,q)=\frac{\partial h(t,q)}{\partial t}\big|_{t=0}$.
\end{lemma}
\begin{proof}
  \sketch
  \leanok
  Define $g(t,q)=\int_0^1\frac{\partial h(ts,q)}{\partial(ts)}\,ds$; after
  changing variables, $t\,g(t,q)=\int_0^t\frac{\partial h(ts,q)}{\partial(ts)}\,d(ts)=h(t,q)$.
  Differentiability of $g$ follows by differentiation under the integral sign
  (the pointwise Fréchet-derivative form of the dominated-convergence theorem, the
  domination bound coming from continuity of $\frac{\partial h}{\partial t}$ on the
  compact set swept by $s\in[0,1]$ and $q$ in a closed ball), and
  $g(0,q)=\int_0^1\frac{\partial h}{\partial t}(0,q)\,ds=\frac{\partial h}{\partial t}(0,q)$.
\end{proof}

A manifold need not be Hausdorff or have a countable coordinate basis. For a
connected manifold these two conditions are equivalent to existence of a
differentiable partition of unity (\cref{thm:dc-ch0-5-6}). Whitney's theorem
gives, for an $n$-manifold with both conditions, an immersion into
$\mathbb{R}^{2n}$ and an embedding into $\mathbb{R}^{2n+1}$; for $n>1$ the
bounds refine to $\mathbb{R}^{2n-1}$ and $\mathbb{R}^{2n}$, respectively.

A family of open sets $V_\alpha\subset M$ with $\bigcup_\alpha V_\alpha=M$ is
\emph{locally finite} if every $p\in M$ has a neighborhood $W$ meeting only
finitely many $V_\alpha$. The \emph{support} of $f:M\to\mathbb{R}$ is the
closure of $\{f\neq0\}$. A family $\{f_\alpha\}$ of differentiable functions is
a \emph{differentiable partition of unity} if $f_\alpha\ge0$, each support lies
in a coordinate neighborhood $V_\alpha$, the family $\{V_\alpha\}$ is locally
finite, and $\sum_\alpha f_\alpha(p)=1$ for every $p\in M$. Such a partition is
\emph{subordinate} to $\{V_\alpha\}$.

\begin{theorem}[existence of partitions of unity]
  \label{thm:dc-ch0-5-6}
  \notready
  \dcref{ch0:5.6}
  A differentiable manifold $M$ has a differentiable partition of unity if and only
  if every connected component of $M$ is Hausdorff and has a countable basis.
\end{theorem}
\begin{proof}
  \uses{lem:dc-ch0-5-6-exists, lem:dc-ch0-5-6-onlyif-complete, lem:dc-ch0-5-6-component}
  \sketch
  For the forward implication, Hausdorffness and a countable basis yield a subordinate
  partition of unity by \cref{lem:dc-ch0-5-6-exists}. Conversely, a partition forces
  Hausdorffness by \cref{lem:dc-ch0-5-6-hausdorff} and second countability by
  \cref{lem:dc-ch0-5-6-onlyif-pou}, using
  \cref{lem:dc-ch0-5-6-sigmacompact-core}. These conclusions combine in
  \cref{lem:dc-ch0-5-6-onlyif-complete}.
\end{proof}

\begin{lemma}[connected space with a locally finite $\sigma$-compact open cover]
  \label{lem:dc-ch0-5-6-sigmacompact-core}
  \lean{Riemannian.sigmaCompactSpace_of_locallyFinite_isSigmaCompact_cover}
  \uses{thm:dc-ch0-5-6}
  \leanok
  \dcref{ch0:5.6}
  Let $X$ be a connected topological space carrying a locally finite open cover
  $\{V_i\}_{i\in\iota}$ each of whose members is $\sigma$-compact. Then $X$ is
  $\sigma$-compact.
\end{lemma}
\begin{proof}
  \sketch
  \leanok
  If every $V_i$ is empty their union is empty, so $X=\varnothing$ is
  $\sigma$-compact; assume then $V_{i_0}\neq\varnothing$ for some $i_0$. Because each
  $V_i=\bigcup_n K_{i,n}$ is a countable union of compacts and the family
  $\{V_j\}$ is locally finite, each compact $K_{i,n}$ meets only finitely many
  $V_j$; hence the overlap set $\{j : V_i\cap V_j\neq\varnothing\}$ is countable for
  every $i$.
  Starting from $\{i_0\}$ and repeatedly adjoining all indices whose $V_j$ meets a
  set already chosen produces, after countably many layers, a countable index set
  $A$ that is closed under overlaps: if $i\in A$ and $V_i\cap V_j\neq\varnothing$
  then $j\in A$. Put $U=\bigcup_{i\in A}V_i$. Then $U$ is open (a union of opens) and
  closed: any $x\in\overline U$ lies in some $V_j$ (the $V$'s cover $X$), and this
  open $V_j$ meets $U$, so $V_j\cap V_i\neq\varnothing$ for some $i\in A$, whence
  $j\in A$ and $x\in V_j\subseteq U$. Thus $U$ is clopen and nonempty, so $U=X$ by
  connectedness. Finally $U=\bigcup_{i\in A}V_i$ is a countable union of
  $\sigma$-compact sets, so $X$ is $\sigma$-compact.
\end{proof}

\begin{lemma}[second countability from a partition of unity]
  \label{lem:dc-ch0-5-6-onlyif}
  \lean{Riemannian.DCSecondCountable_of_locallyFinite_isSigmaCompact_cover}
  \uses{lem:dc-ch0-5-6-sigmacompact-core, thm:dc-ch0-5-6}
  \leanok
  \dcref{ch0:5.6}
  Let $M$ be a connected manifold charted over a second-countable model $H$. If $M$
  carries a locally finite open cover by $\sigma$-compact sets — in particular, the
  coordinate neighborhoods $\{V_\alpha\}$ of a differentiable partition of unity,
  which are locally finite, cover $M$, and are each homeomorphic to an open subset
  of $\mathbb{R}^n$ (hence $\sigma$-compact) — then $M$ has a countable basis.
\end{lemma}
\begin{proof}
  \sketch
  \leanok
  By \cref{lem:dc-ch0-5-6-sigmacompact-core} the manifold $M$ is $\sigma$-compact.
  A charted space over a second-countable model that is $\sigma$-compact is itself
  second countable (cover each of the countably many compacts by finitely many chart
  domains and use the countable basis of $H$ inside each). Applying this argument to
  each connected component gives the second-countability
  conclusion in the theorem.
\end{proof}

\begin{lemma}[second countability from an honest partition of unity]
  \label{lem:dc-ch0-5-6-onlyif-pou}
  \lean{Riemannian.DCSecondCountable_of_partitionOfUnity}
  \uses{lem:dc-ch0-5-6-onlyif, thm:dc-ch0-5-6}
  \leanok
  \dcref{ch0:5.6}
  Let $M$ be a connected manifold charted over a second-countable model $H$, and let
  $\{f_\alpha\}$ be a differentiable partition of unity on $M$ (with $\sum_\alpha f_\alpha
  \equiv 1$) subordinate to a locally finite family $\{V_\alpha\}$ of $\sigma$-compact
  coordinate neighborhoods. Then $M$ has a countable basis.
\end{lemma}
\begin{proof}
  \sketch
  \leanok
  Unlike \cref{lem:dc-ch0-5-6-onlyif}, the covering hypothesis is not assumed but
  \emph{derived} from the partition. For each $p\in M$, condition $\sum_\alpha f_\alpha(p)=1$
  gives some index $\alpha$ with $f_\alpha(p)\neq0$, so $p\in\operatorname{supp}f_\alpha
  \subseteq\operatorname{tsupp}f_\alpha\subseteq V_\alpha$ by subordinacy; hence
  $\{V_\alpha\}$ covers $M$. As these neighborhoods are open, $\sigma$-compact, and locally
  finite, \cref{lem:dc-ch0-5-6-onlyif} (equivalently
  \cref{lem:dc-ch0-5-6-sigmacompact-core} followed by charted second-countability) yields a
  countable basis for $M$.
\end{proof}

\begin{lemma}[a partition of unity forces the Hausdorff axiom]
  \label{lem:dc-ch0-5-6-hausdorff}
  \lean{Riemannian.t2Space_of_partitionOfUnity}
  \uses{thm:dc-ch0-5-6}
  \leanok
  \dcref{ch0:5.6}
  Let $M$ be a manifold admitting a differentiable partition of unity $\{f_\alpha\}$
  subordinate to a family $\{V_\alpha\}$ of open, Hausdorff coordinate neighborhoods. Then
  $M$ is Hausdorff.
\end{lemma}
\begin{proof}
  \sketch
  \leanok
  Let $p\neq q$ have no disjoint neighborhoods; then the filter $\mathcal F=\mathcal N_p
  \sqcap\mathcal N_q$ is proper. Each $f_\alpha$ is continuous, so it tends to both
  $f_\alpha(p)$ and $f_\alpha(q)$ along $\mathcal F$; uniqueness of limits in $\mathbb R$
  forces $f_\alpha(p)=f_\alpha(q)$. Since $\sum_\alpha f_\alpha(p)=1$, some $f_\alpha(p)\neq
  0$, whence $f_\alpha(q)\neq0$ too, so both $p,q\in\operatorname{supp}f_\alpha\subseteq
  V_\alpha$. But $V_\alpha$ is open and Hausdorff, so $p,q$ are separated by sets open in
  $V_\alpha$ and hence open in $M$ — contradicting the choice of $p,q$. Thus $M$ is
  Hausdorff.
\end{proof}

\begin{lemma}[complete ``only if'' direction from a partition of unity]
  \label{lem:dc-ch0-5-6-onlyif-complete}
  \lean{Riemannian.DCManifoldAxioms_of_partitionOfUnity}
  \uses{lem:dc-ch0-5-6-hausdorff, lem:dc-ch0-5-6-onlyif-pou}
  \leanok
  \dcref{ch0:5.6}
  Let $M$ be a connected manifold charted over a second-countable model $H$ admitting a
  differentiable partition of unity subordinate to a locally finite family of open,
  Hausdorff, $\sigma$-compact coordinate neighborhoods. Then $M$ is Hausdorff and has a
  countable basis.
\end{lemma}
\begin{proof}
  \sketch
  \leanok
  Immediate from \cref{lem:dc-ch0-5-6-hausdorff} (Hausdorff) and
  \cref{lem:dc-ch0-5-6-onlyif-pou} (countable basis), component by component.
\end{proof}

\begin{lemma}[existence of a subordinate partition of unity]
  \label{lem:dc-ch0-5-6-exists}
  \lean{Riemannian.DCSmoothPartitionOfUnity.exists_isSubordinate}
  \uses{thm:dc-ch0-5-6}
  \leanok
  \dcref{ch0:5.6}
  Let $M$ be a Hausdorff, second-countable (equivalently $\sigma$-compact),
  finite-dimensional differentiable manifold. For every closed set $s\subseteq M$
  and every open covering $\{U_i\}$ of $s$ there is a differentiable partition of
  unity $\{f_i\}$ subordinate to $\{U_i\}$, i.e. with
  $\operatorname{supp} f_i\subseteq U_i$.
\end{lemma}
\begin{proof}
  \sketch
  \leanok
  Hausdorffness and second countability make $M$ paracompact and $\sigma$-compact.
  Local Euclidean bump functions supported in the $U_i$ form a locally finite family;
  normalizing their sum yields the required partition.
\end{proof}

\begin{lemma}[normalization of a partition of unity]
  \label{lem:dc-ch0-5-6-sum}
  \lean{Riemannian.DCSmoothPartitionOfUnity.sum_eq_one}
  \leanok
  \dcref{ch0:5.6}
  Let $\{f_\alpha\}$ be a differentiable partition of unity on $M$ subordinate to a
  covering of a set $s\subseteq M$. Then $\sum_\alpha f_\alpha(p)=1$ for every
  $p\in s$.
\end{lemma}
\begin{proof}
  \sketch
  \leanok
  This is condition (3) in the definition of a differentiable partition of unity:
  the family is locally finite, so the sum is finite at each point, and by
  definition it takes the value $1$ throughout the covered set $s$.
\end{proof}

\begin{remark}[bump functions and locality]
  \label{rem:dc-ch0-5-7}
  \lean{Riemannian.exists_ch0_bump_function, Riemannian.exists_ch0_manifold_bump_function}
  \leanok
  \dcref{ch0:5.7}
  Given $p\in\mathbb{R}^n$ and an open ball $B_r(p)\subset\mathbb{R}^n$, there exist
  a neighborhood $U$ of $p$ with $\overline U\subset B_r(p)$ and a differentiable
  $f:\mathbb{R}^n\to\mathbb{R}$ with $0\le f(q)\le1$ and
  $$
  f(q)=\begin{cases}1,&q\in\overline U,\\ 0,&q\notin B_r(p).\end{cases}
  $$
  For $r=2$ take $U=B_1(p)$ and $f(q)=\beta(-|p-q|)$, where
  $\beta(t)=\frac{\int_{-\infty}^t\alpha(s)\,ds}{\int_{-2}^{-1}\alpha(s)\,ds}$ and
  $\alpha:\mathbb{R}\to\mathbb{R}$ equals $\exp\big(-\frac{1}{(t+2)(-1-t)}\big)$ for
  $-2<t<-1$ and zero otherwise. The same holds on $M$: if $V\subset M$ is a
  neighborhood of $p$ contained in a coordinate neighborhood diffeomorphic to an
  open ball, there exist $U$ with $\overline U\subset V$ and differentiable
  $f:M\to\mathbb{R}$, $0\le f\le1$, $f=1$ on $\overline U$, $f=0$ off $V$.
  These bump functions localize globally defined constructions.
\end{remark}

\begin{example}[a non-differentiable curve]
  \label{ex:dc-ch0-3-2}
  \dcref{ch0:3.2}
  \lean{Riemannian.ch0AbsCurve, Riemannian.ch0AbsCurve_not_differentiableAt_zero}
  \leanok
  The curve $\alpha:\mathbb{R}\to\mathbb{R}^2$, $\alpha(t)=(t,|t|)$, is not
  differentiable at $t=0$ (Fig. 5).
\end{example}

\begin{example}[differentiable but not an immersion]
  \label{ex:dc-ch0-3-3}
  \dcref{ch0:3.3}
  \lean{Riemannian.ch0CuspCurve, Riemannian.ch0CuspCurve_differentiable,
    Riemannian.ch0CuspCurve_fderiv_zero, Riemannian.ch0CuspCurve_not_immersion}
  \leanok
  $\alpha:\mathbb{R}\to\mathbb{R}^2$, $\alpha(t)=(t^3,t^2)$, is differentiable but
  not an immersion: the immersion condition is $\alpha'(t)\neq0$, which fails at
  $t=0$ (Fig. 6).
\end{example}

\begin{example}[immersion with self-intersection]
  \label{ex:dc-ch0-3-4}
  \dcref{ch0:3.4}
  \lean{Riemannian.ch0SelfIntersectionCurve,
    Riemannian.ch0SelfIntersectionCurve_hasDerivAt,
    Riemannian.ch0SelfIntersectionCurve_differentiable,
    Riemannian.ch0SelfIntersectionCurve_deriv_ne_zero,
    Riemannian.ch0SelfIntersectionCurve_self_intersects}
  \leanok
  $\alpha(t)=(t^3-4t,\,t^2-4)$ (Fig. 7) is an immersion
  $\alpha:\mathbb{R}\to\mathbb{R}^2$ with a self-intersection at $t=2,t=-2$. Hence
  $\alpha$ is not an embedding.
\end{example}

\begin{example}[injective immersion that is not an embedding]
  \label{ex:dc-ch0-3-5}
  \dcref{ch0:3.5}
  \notready
  The curve $\alpha:(-3,0)\to\mathbb{R}^2$ (Fig. 8), defined piecewise as the
  segment $(0,-(t+2))$ for $t\in(-3,-1)$, a regular joining curve for
  $t\in(-1,-\frac1\pi)$, and $(-t,-\sin\frac1t)$ for $t\in(-\frac1\pi,0)$, is an
  immersion without self-intersections that is not an embedding: a neighborhood of
  a point $p$ in the vertical part has infinitely many connected components in the
  topology induced from $\mathbb{R}^2$, whereas in the topology induced from
  $\alpha$ (the line) it is an open interval, hence connected.
\end{example}

\begin{example}[the sphere $S^n$ is orientable]
  \label{ex:dc-ch0-4-6}
  \dcref{ch0:4.6}
  \notready
  \uses{ex:dc-ch0-4-5}
  $S^n=\{(x_1,\dots,x_{n+1})\in\mathbb{R}^{n+1};\sum_{i=1}^{n+1}x_i^2=1\}$. With
  north/south poles $N=(0,\dots,0,1)$, $S=(0,\dots,0,-1)$, the *stereographic
  projections* $\pi_1:S^n-\{N\}\to\mathbb{R}^n$,
  $\pi_1(x)=\big(\frac{x_1}{1-x_{n+1}},\dots,\frac{x_n}{1-x_{n+1}}\big)$, and
  $\pi_2:S^n-\{S\}\to\mathbb{R}^n$ (from the south pole) are differentiable
  injections onto $x_{n+1}=0$. The parametrizations
  $(\mathbb{R}^n,\pi_1^{-1}),(\mathbb{R}^n,\pi_2^{-1})$ cover $S^n$ with change of
  coordinates $y_j'=\frac{y_j}{\sum_{i=1}^n y_i^2}$, differentiable. For $n\geq2$,
  $\pi_1^{-1}(\mathbb{R}^n)\cap\pi_2^{-1}(\mathbb{R}^n)=S^n-\{N,S\}$ is connected,
  so by \cref{ex:dc-ch0-4-5}, $S^n$ is orientable. For $n=1$, the same atlas is
  orientable directly: on the two components of $S^1-\{N,S\}$, the change of
  coordinates is $y'=1/y$ with derivative $-1/y^2<0$, and changing the sign of one
  coordinate makes both transition derivatives positive. The antipodal map
  $A:S^n\to S^n$, $A(p)=-p$, is a diffeomorphism ($A^2=\mathrm{id}$); it reverses the
  orientation of $S^n$ when $n$ is even and preserves it when $n$ is odd.
\end{example}

\begin{example}[special cases of Example 4.8]
  \label{ex:dc-ch0-4-9}
  \dcref{ch0:4.9}
  \notready
  \uses{ex:dc-ch0-4-8}
  **4.9 (a) (the $k$-torus).** Let $G$ be the group of "integral" translations of
  $\mathbb{R}^k$: for $(n_1,\dots,n_k)\in\mathbb{Z}^k$, the corresponding
  translation sends $(x_1,\dots,x_k)$ to $(x_1+n_1,\dots,x_k+n_k)$.
  This is a properly discontinuous action; $\mathbb{R}^k/G$ (with the structure of
  \cref{ex:dc-ch0-4-8}) is the *$k$-torus $T^k$*. For $k=2$, $T^2$ is diffeomorphic to the
  torus of revolution in $\mathbb{R}^3$, the inverse image of $0$ of
  $f(x,y,z)=z^2+(\sqrt{x^2+y^2}-a)^2-r^2$.
  \textbf{4.9 (b) (Klein bottle; Möbius band).} Let $S\subset\mathbb{R}^3$ be a regular
  surface symmetric about $0$ (if $p\in S$ then $-p=A(p)\in S$). The group
  $\{A,\mathrm{Id}\}$ acts properly discontinuously on $S$; $S/G$ carries the
  structure of \cref{ex:dc-ch0-4-8}. When $S$ is the torus of revolution $T^2$, $S/G=K$ is
  \emph{the Klein bottle}; when $S=C=\{(x,y,z); x^2+y^2=1,\,-1<z<1\}$ (a right circular
  cylinder), $S/G$ is the \emph{Möbius band}. By Exercise 9 the Klein bottle and Möbius
  band are non-orientable.
\end{example}

\begin{lemma}[fixed-cover characterization on an open connected component]
  \label{lem:dc-ch0-5-6-component}
  \dcref{ch0:5.6}
  \uses{lem:dc-ch0-5-6-fixed-cover}
  \lean{Riemannian.DCPartitionOfUnity_iff_manifoldAxioms_on_connectedComponent}
  \leanok
  Let $p\in M$ and suppose its connected component $C_p$ is open in $M$. For a locally finite
  cover of $C_p$ by open Hausdorff $\sigma$-compact sets, a smooth partition of unity on $C_p$
  subordinate to this fixed cover exists if and only if $C_p$ is Hausdorff and has a countable
  basis. The cover and partition are genuinely component-local; no global countability or
  cross-component assembly is asserted.
\end{lemma}
\begin{proof}
  \leanok
  Regard $C_p$ as the open subspace of $M$. Connected components are preconnected, so this
  subspace carries the connectedness instance required by
  \cref{lem:dc-ch0-5-6-fixed-cover}; applying that characterization to the inherited manifold
  structure gives the result.
\end{proof}

\begin{lemma}[fixed-cover partition-of-unity characterization]
  \label{lem:dc-ch0-5-6-fixed-cover}
  \dcref{ch0:5.6}
  \lean{Riemannian.DCPartitionOfUnity_iff_manifoldAxioms}
  \leanok
  Let $M$ be a connected finite-dimensional manifold, charted over a second-countable model,
  and let $\{V_i\}_{i\in\iota}$ be a locally finite cover by open Hausdorff $\sigma$-compact
  sets. Then a smooth partition of unity subordinate to this fixed cover exists if and only if
  $M$ is Hausdorff and has a countable basis.
\end{lemma}
\begin{proof}
  \leanok
  This is the connected fixed-cover form of the theorem: the forward implication obtains the
  Hausdorff and second-countability axioms from the partition, while the reverse implication
  applies the smooth subordinate-partition construction to the closed set $M$ itself.
\end{proof}
